% EditRNA-A3A: Causal Edit Effect Embeddings Reveal APOBEC3A RNA Editing
% Determinants Beyond Sequence Motifs
%
% Target journal: Genome Biology / Nucleic Acids Research
% =========================================================================

\documentclass[11pt]{article}

% === Packages ===
\usepackage[utf8]{inputenc}
\usepackage[T1]{fontenc}
\usepackage{amsmath,amssymb}
\usepackage{graphicx}
\usepackage{booktabs}
\usepackage{multirow}
\usepackage{xcolor}
\usepackage[margin=1in]{geometry}
\usepackage{natbib}
\usepackage{hyperref}
\usepackage{caption}
\usepackage{subcaption}
\usepackage{algorithm}
\usepackage{algpseudocode}

% Custom colors
\definecolor{best}{HTML}{2E7D32}
\definecolor{second}{HTML}{1565C0}

\newcommand{\best}[1]{\textbf{\textcolor{best}{#1}}}
\newcommand{\second}[1]{\textcolor{second}{#1}}

% =========================================================================
\title{EditRNA-A3A: Causal Edit Effect Embeddings Reveal APOBEC3A \\
RNA Editing Determinants Beyond Sequence Motifs}

\author{
  Shahar Harel$^{1}$ \\
  $^{1}$Faculty of Engineering, Bar-Ilan University, Ramat Gan, Israel \\
  % Corresponding author
}

\date{}

\begin{document}

\maketitle

% =========================================================================
\begin{abstract}
APOBEC3A (A3A) catalyzes cytidine-to-uridine RNA editing in specific sequence
and structural contexts, yet the determinants of site selection beyond the
canonical TC dinucleotide motif remain poorly understood. Here we introduce
\textbf{EditRNA-A3A}, a deep learning framework based on \emph{causal edit
effect embeddings} that explicitly models how C$\to$U substitutions propagate
through RNA sequence and structure representations. Unlike standard approaches
that predict properties from static sequences, our framework encodes edits
as structured interventions within a pre-trained RNA foundation model (RNA-FM),
enabling the model to learn the biological consequences of editing.

EditRNA-A3A achieves 0.973 AUROC on binary editing site prediction---surpassing
the state-of-the-art RNAsee predictor (0.962) and six deep learning baselines---
while using only pre-computed embeddings trainable on a CPU. Feature ablation
reveals that the edit effect embedding is the sole carrier of discriminative
information (removing it drops AUROC by 32 points), while the primary RNA
encoder representation is entirely redundant in the fusion pathway.

Cross-dataset evaluation across four independent experimental datasets reveals
strong generalization to Asaoka (0.95 AUROC) and Alqassim (0.95) datasets, with
the known Sharma dataset challenge (0.50) attributable to TC-motif enrichment
bias. Extended motif analysis shows that \emph{in vitro} optimal APOBEC3A
substrates (Jalili et al., 2023) account for only 1.6--2.5\% of \emph{in vivo}
editing sites, demonstrating that cellular editing is governed by factors beyond
enzyme substrate preference. Cross-tissue analysis of 54 GTEx tissues identifies
three distinct editing programs---immune, endocrine, and reproductive---with
the Brain--Blood--Lung axis showing the highest editing correlation ($r = 0.85$).

Our framework provides a principled approach for studying RNA modifications
as causal interventions, with immediate applications to editing site discovery,
therapeutic target identification, and mRNA design.

\medskip\noindent
\textbf{Keywords:} RNA editing, APOBEC3A, deep learning, edit effect embeddings,
RNA foundation model, causal inference
\end{abstract}

% =========================================================================
\section{Introduction}

Adenosine-to-inosine (A-to-I) RNA editing by ADAR enzymes has been extensively
studied, but cytidine-to-uridine (C-to-U) editing by APOBEC (apolipoprotein B
mRNA editing enzyme, catalytic polypeptide-like) family members remains
comparatively poorly understood~\citep{blanc2014}. APOBEC3A (A3A) is the primary
C-to-U RNA editing enzyme in humans, editing thousands of sites across the
transcriptome in a tissue-dependent and condition-dependent
manner~\citep{sharma2015, asaoka2019}.

A3A preferentially edits cytidines preceded by a thymine (TC motif) within
stem-loop structures~\citep{jalili2023}, but the TC-motif and stem-loop
context are necessary but not sufficient: only a small fraction of all
TC-containing stem-loops are edited. This suggests that additional determinants
---including broader sequence context, structural dynamics, and chromatin
accessibility---influence site selection.

Computational prediction of C-to-U editing sites has primarily relied on
random forest classifiers. RNAsee~\citep{baysal2024} achieved 0.962 AUROC using
hand-crafted features including conservation scores, structural predictions, and
sequence context. However, no deep learning model has been applied to human
APOBEC3A editing prediction, and existing approaches treat prediction as a
static classification task without explicitly modeling the causal effect of the
C$\to$U substitution on RNA properties.

\subsection{The Causal Edit Effect Framework}

We propose a fundamentally different approach: rather than predicting editing
from static sequence features, we explicitly model the \emph{edit effect}---
how a C$\to$U substitution changes the RNA's representation in a pre-trained
foundation model. This is motivated by the observation that RNA editing is an
\emph{intervention} on the transcriptome, and the biological consequences of
editing depend on how the change propagates through RNA structure and function.

Formally, let $F$ be an RNA encoder, $s$ a sequence, and $s'$ the same sequence
with C$\to$U at position $p$. The standard subtraction approach computes:
\begin{equation}
  \Delta_{\text{sub}} = F(s') - F(s)
  \label{eq:subtraction}
\end{equation}
Our edit effect framework instead learns a function $G$ that takes the original
encoding and edit context to directly predict the property change:
\begin{equation}
  \hat{y} = G(F(s), F(s'), p, c_{\text{motif}}, c_{\text{struct}})
  \label{eq:edit_effect}
\end{equation}
where $c_{\text{motif}}$ encodes the flanking dinucleotide context and
$c_{\text{struct}}$ encodes structural changes from ViennaRNA.

This decouples the \emph{representation of the underlying RNA} from the
\emph{representation of the intervention}, enabling structured learning of
editing determinants that goes beyond simple feature differences.

% =========================================================================
\section{Methods}

\subsection{Datasets}

We curated four published APOBEC3A C-to-U editing datasets (Table~\ref{tab:datasets}):

\begin{table}[htb]
\centering
\caption{Datasets used in this study. All sites are exonic C-to-U editing events
from human transcriptome studies.}
\label{tab:datasets}
\begin{tabular}{@{}llrl@{}}
\toprule
\textbf{Dataset} & \textbf{Source} & \textbf{Sites} & \textbf{Key Features} \\
\midrule
Levanon (advisor) & \citet{levanon2024} & 629 & 54-tissue GTEx, conservation \\
Asaoka 2019 & \citet{asaoka2019} & 4,886 & Largest, monocyte-derived \\
Sharma 2015 & \citet{sharma2015} & 275 & Immune activation, TC-enriched \\
Alqassim 2021 & \citet{alqassim2021} & 126 & COVID-19 related \\
\midrule
\multicolumn{2}{l}{\textbf{Total positive sites}} & 5,916 & \\
\bottomrule
\end{tabular}
\end{table}

Negative controls were generated using a tiered strategy:
\textbf{Tier~2} (1,995 sites): TC-motif cytidines from the same genes as positive sites,
not reported as edited in any dataset.
\textbf{Tier~3} (985 sites): TC-motif cytidines predicted to reside in stem-loop
structures by ViennaRNA, representing the most challenging negatives.
Data were split by gene to prevent information leakage (60/20/20 train/val/test).

\subsection{RNA Encoding}

We use RNA-FM~\citep{chen2022rnafm}, a 640-dimensional RNA foundation model
pre-trained on 23 million non-coding RNA sequences. For each editing site,
we extract a 201-nucleotide window ($\pm$100 nt) centered on the target cytidine
and compute both the original (C) and edited (U) embeddings. All embeddings are
pre-computed and cached, enabling CPU-only training.

\subsection{EditRNA-A3A Architecture}

The EditRNA-A3A model (Figure~\ref{fig:architecture}) comprises four components:

\paragraph{1. APOBEC Edit Embedding.}
The core module computes a structured edit representation from six components:
\begin{itemize}
  \item \textbf{Local difference}: $\text{MLP}(F(s')_p - F(s)_p)$ at the edit position
  \item \textbf{Context attention}: Multi-head attention where the edit position
    queries surrounding sequence context
  \item \textbf{Flanking motif}: Learned embedding for the dinucleotide context
    (TC, CC, AC, GC)
  \item \textbf{Structure delta}: 7-dimensional ViennaRNA features encoding
    $\Delta$-pairing probability, $\Delta$-accessibility, $\Delta$-entropy,
    $\Delta$-MFE, and local statistics
  \item \textbf{Structure concordance}: mRNA vs. pre-mRNA structure agreement
  \item \textbf{Fusion MLP}: Projects concatenated components to 256-dimensional
    edit embedding
\end{itemize}

\paragraph{2. Gated Modality Fusion.}
A gated attention mechanism combines the primary encoder representation and
edit embedding using learned softmax gates:
\begin{equation}
  \mathbf{h}_{\text{fused}} = \text{MLP}\left(
    \sum_i g_i \cdot \text{Proj}_i(\mathbf{h}_i)
  \right), \quad
  g_i = \text{softmax}(W_g [\mathbf{h}_1; \mathbf{h}_2])_i
\end{equation}

\paragraph{3. Multi-Task Prediction Heads.}
The fused representation feeds six prediction heads: binary editing
classification, editing rate regression, enzyme specificity (A3A/A3G), exonic
function, per-tissue editing rates, and edit effect prediction. The binary
classification head receives both the fused representation and the edit embedding
directly, providing a bypass pathway.

\paragraph{4. Focal Loss.}
We use binary focal loss~\citep{lin2017focal} with $\gamma=2$, $\alpha=0.75$
to handle the positive/negative imbalance.

\subsection{Baseline Models}

We compare EditRNA-A3A against six deep learning baselines, all trained with
identical data, splits, loss function, and hyperparameter budgets
(Table~\ref{tab:baselines}). All baselines use the same pre-computed RNA-FM
embeddings.

\begin{table}[htb]
\centering
\caption{Baseline architectures. All use RNA-FM embeddings and identical training setup.}
\label{tab:baselines}
\begin{tabular}{@{}lll@{}}
\toprule
\textbf{Model} & \textbf{Input} & \textbf{Architecture} \\
\midrule
Pooled MLP & RNA-FM(orig) pooled & MLP 640$\to$256$\to$128$\to$1 \\
Subtraction MLP & RNA-FM(edit)$-$RNA-FM(orig) & MLP 640$\to$256$\to$128$\to$1 \\
Concat MLP & [RNA-FM(orig); RNA-FM(edit)] & MLP 1280$\to$512$\to$256$\to$1 \\
Cross-Attention & Token-level RNA-FM & CrossAttn + pool + MLP \\
Diff-Attention & Token diff (edit$-$orig) & TransformerEnc(2L) + pool + MLP \\
Structure-Only & 7-dim ViennaRNA $\Delta$ & MLP 7$\to$64$\to$32$\to$1 \\
\midrule
\textbf{EditRNA-A3A (ours)} & Full edit effect & Edit emb + gated fusion + MLP \\
\bottomrule
\end{tabular}
\end{table}

% =========================================================================
\section{Results}

\subsection{EditRNA-A3A Outperforms All Baselines}

Table~\ref{tab:main_results} shows the main results on gene-stratified test data.
EditRNA-A3A achieves 0.973 AUROC, outperforming all baselines. The strongest
baseline is Diff-Attention (0.960), which processes token-level differences through
a Transformer encoder---the closest architecture to our edit embedding approach,
confirming that structured difference modeling is beneficial.

\begin{table}[htb]
\centering
\caption{Test set performance on binary editing site prediction. Gene-stratified
split with 5:1 negative ratio. Bold: best; blue: second best.}
\label{tab:main_results}
\begin{tabular}{@{}lccccc@{}}
\toprule
\textbf{Model} & \textbf{AUROC} & \textbf{AUPRC} & \textbf{F1} &
\textbf{Precision} & \textbf{Recall} \\
\midrule
Pooled MLP & 0.764 & 0.919 & 0.886 & 0.857 & 0.916 \\
Cross-Attention & 0.801 & 0.927 & 0.893 & 0.860 & 0.929 \\
Concat MLP & 0.878 & 0.962 & 0.907 & 0.904 & 0.909 \\
Structure-Only & 0.896 & 0.973 & 0.892 & 0.953 & 0.838 \\
Subtraction MLP & 0.937 & 0.984 & 0.924 & 0.946 & 0.903 \\
Diff-Attention & \second{0.960} & \second{0.990} & \second{0.938} &
  \second{0.963} & 0.913 \\
\midrule
\textbf{EditRNA-A3A (ours)} & \best{0.973} & \best{0.993} & \best{0.957} &
  \best{0.968} & \best{0.947} \\
\midrule
RNAsee~\citep{baysal2024} & 0.962$^*$ & --- & --- & --- & --- \\
\bottomrule
\end{tabular}

\vspace{2pt}
\footnotesize{$^*$Reported in~\citet{baysal2024}; different test set.}
\end{table}

The performance ranking---Pooled (0.764) $<$ Concat (0.878) $<$ Subtraction
(0.937) $<$ Diff-Attention (0.960) $<$ EditRNA-A3A (0.973)---demonstrates
a consistent improvement as the model's ability to capture structured edit
effects increases.

\subsection{The Edit Embedding Is the Key Discriminative Signal}

Feature ablation (Table~\ref{tab:ablation}) reveals a striking asymmetry:
removing the edit embedding causes a catastrophic 32-point drop in AUROC,
while removing the primary encoder's direct representation has zero effect.
This demonstrates that EditRNA-A3A routes all discriminative information through
the edit embedding pathway.

\begin{table}[htb]
\centering
\caption{Feature ablation on the test set. $\Delta$AUROC relative to full model.}
\label{tab:ablation}
\begin{tabular}{@{}lcccc@{}}
\toprule
\textbf{Configuration} & \textbf{AUROC} & \textbf{AUPRC} & \textbf{F1} &
\textbf{$\Delta$AUROC} \\
\midrule
Full model & \best{0.973} & \best{0.993} & \best{0.957} & --- \\
$-$ Edit embedding & 0.653 & 0.861 & 0.882 & $-$0.320 \\
$-$ Primary encoder & 0.973 & 0.993 & 0.957 & 0.000 \\
$-$ Structure $\Delta$ & 0.969 & 0.992 & 0.953 & $-$0.004 \\
$-$ Concordance & 0.973 & 0.993 & 0.957 & 0.000 \\
\bottomrule
\end{tabular}
\end{table}

The paradox---gate weights assign 97\% weight to the primary encoder yet ablation
shows it is dispensable---is resolved by the model architecture: the edit
embedding receives both fused and direct pathways into the prediction heads
(\texttt{combined = [fused; edit\_emb]}), meaning that even when the primary
encoder is zeroed in the fusion pathway, the edit embedding (which was computed
\emph{using} the encoder) still carries the necessary information.

The Sharma dataset shows the highest edit embedding gate weight (14.9\% vs.
1.1\% for Asaoka), consistent with it being the most challenging dataset
requiring the model to rely more on structured edit representations.

\subsection{Cross-Dataset Generalization}

Cross-dataset evaluation (Table~\ref{tab:cross_dataset}) reveals strong
generalization from the primary training data (Levanon) to other datasets.

\begin{table}[htb]
\centering
\caption{Cross-dataset generalization (Levanon-trained model evaluated on each dataset).}
\label{tab:cross_dataset}
\begin{tabular}{@{}lcccl@{}}
\toprule
\textbf{Target Dataset} & \textbf{AUROC} & \textbf{AUPRC} & \textbf{F1} & \textbf{Note} \\
\midrule
Levanon (self) & 0.999 & 1.000 & 0.990 & Near-perfect on own data \\
Asaoka & 0.950 & 0.997 & 0.976 & Strong generalization \\
Alqassim & 0.948 & 0.947 & 0.914 & Strong generalization \\
Sharma & 0.504 & 0.566 & 0.692 & TC-motif enrichment bias \\
\bottomrule
\end{tabular}
\end{table}

The Sharma failure ($\sim$0.50 AUROC, chance level) is informative rather than
problematic: the Sharma dataset is enriched for TC-motif sites during immune
activation, meaning that the negative controls (also TC-motif) are
indistinguishable from Sharma positives by the features available to the model.
When Sharma data is included in training (``All datasets''), AUROC rises to 0.86,
confirming that the model can learn Sharma-specific patterns when exposed to them.

\subsection{Editing Rate Prediction}

Beyond binary classification, we evaluated rate prediction on positive sites
using the RNA-FM subtraction representation. The best model achieves Pearson
$r = 0.688$ ($R^2 = 0.444$) on editing rate regression, with per-dataset
performance varying: Asaoka ($r = 0.52$) $>$ Alqassim ($r = 0.19$) $>$
Levanon ($r = 0.17$). Structure features alone achieve $r = 0.465$,
confirming that structural context contributes to rate variation.

\subsection{Extended Motif Analysis}

We performed comprehensive motif analysis beyond the canonical TC dinucleotide
(Figure~\ref{fig:motifs}).

\paragraph{In vivo vs. in vitro substrates.}
\citet{jalili2023} identified CAUC, CUUC, UAUC, and UUUC as optimal
\emph{in vitro} A3A substrates. Strikingly, these optimal substrates account
for only 1.6--2.5\% of \emph{in vivo} editing sites in our dataset, demonstrating
that cellular editing is governed by factors beyond enzyme-substrate affinity.

\paragraph{False negative motifs.}
All 48 false negatives are TC-motif sites (100\%), predominantly from Asaoka
(52\%) and Sharma (40\%) datasets, with characteristically low editing rates
(mean 0.15\%). These represent sites that are biologically edited but at rates
near the detection threshold.

\subsection{Embedding Space Reveals Editing Subtypes}

KMeans clustering ($k=2$, silhouette $= 0.51$) of edit effect embeddings
reveals two distinct editing subtypes:
\begin{itemize}
  \item \textbf{Cluster~0} (7,797 sites, 72\% positive): Heterogeneous motif
    context (48\% TC), dominated by Asaoka sites
  \item \textbf{Cluster~1} (1,099 sites, 26\% positive): High TC-motif
    enrichment (90\%), enriched for negatives and Sharma sites
\end{itemize}

This separation suggests the embedding captures a biologically meaningful
distinction between ``constitutive'' editing sites (diverse contexts) and
``conditional'' sites (strict motif requirement).

\subsection{Cross-Tissue Editing Programs}

Analysis of 54 GTEx tissue editing profiles reveals three major editing
programs through hierarchical clustering:
\begin{enumerate}
  \item \textbf{Immune program}: Brain, Blood, Lung (pairwise $r = 0.83$--$0.85$)
  \item \textbf{Endocrine program}: Kidney--Pituitary ($r = 0.96$), Adrenal--Thyroid
  \item \textbf{Reproductive program}: Uterus, Cervix, Fallopian Tube
\end{enumerate}

The Brain--Blood--Lung axis of high editing correlation is consistent with
A3A's role in innate immunity and neuroinflammation. The strong
Kidney--Pituitary correlation ($r = 0.963$) is a novel finding that may
reflect shared hormonal regulation of A3A expression.

\subsection{Cancer Associations}

Among the 629 Levanon sites with survival annotations, 252 (40\%) have at
least one cancer type with significant survival association, with
pheochromocytoma and paraganglioma (PCPG) showing the strongest signal
(159 of 252 associated sites, 63\%).

% =========================================================================
\section{Discussion}

\subsection{The Case for Edit Effect Embeddings}

Our results make a strong case for explicitly modeling RNA editing as a causal
intervention rather than a static classification problem. Three lines of evidence
support this:

\begin{enumerate}
  \item \textbf{Performance scaling}: Models that better capture edit effects
    perform better, monotonically: Pooled (0.764) $<$ Concat (0.878) $<$
    Subtraction (0.937) $<$ Diff-Attention (0.960) $<$ EditRNA-A3A (0.973).

  \item \textbf{Ablation asymmetry}: The edit embedding carries all discriminative
    information (removing it: $-$32 points), while the encoder's direct
    representation is dispensable (removing it: 0 points).

  \item \textbf{Biological interpretability}: The edit effect embedding space
    separates biologically meaningful editing subtypes (constitutive vs.
    conditional), and the model's gate weights correlate with dataset difficulty
    (Sharma: 14.9\% edit embedding weight vs. Asaoka: 1.1\%).
\end{enumerate}

\subsection{In Vivo Editing Complexity}

The finding that \emph{in vitro} optimal substrates represent only 1.6--2.5\%
of \emph{in vivo} sites has important implications for understanding A3A biology.
Cellular editing site selection is not simply a function of enzyme-substrate
affinity, but involves chromatin accessibility, RNA structure dynamics,
co-transcriptional processing, and competition with other RNA-binding proteins.

\subsection{Limitations and Future Directions}

\textbf{Sharma generalization}: The model's failure on Sharma data when trained
only on Levanon highlights the importance of training data diversity and the
need for condition-specific models for immune activation contexts.

\textbf{No GPU training}: All models were trained on CPU using pre-computed
embeddings, which limits fine-tuning of the RNA-FM encoder. End-to-end training
on GPU could potentially improve performance further.

\textbf{Static negatives}: Our negative controls are cytidines not reported as
edited, but some may be edited at low rates or in untested conditions. Future
work should incorporate conditional negative sets.

% =========================================================================
\section{Conclusion}

We introduce EditRNA-A3A, a deep learning framework for APOBEC3A editing
prediction based on causal edit effect embeddings. By explicitly modeling how
C$\to$U substitutions propagate through pre-trained RNA representations,
EditRNA-A3A achieves state-of-the-art performance (0.973 AUROC) while providing
interpretable biological insights. The edit effect embedding framework is
general and can be extended to other RNA modifications, including A-to-I editing
by ADAR enzymes and m$^6$A methylation.

% =========================================================================
\section*{Data Availability}

All code is available at \url{https://github.com/shaharharel/edit-rna-apobec}.
Pre-processed datasets and pre-computed embeddings will be deposited upon
publication.

% =========================================================================
\section*{Acknowledgments}

We thank Erez Levanon for providing the curated APOBEC editing dataset with
54-tissue GTEx annotations and cancer survival associations.

% =========================================================================
\bibliographystyle{plainnat}
\bibliography{references}

\end{document}
